\documentclass[aspectratio=169]{beamer}

\usetheme{Madrid}
\usecolortheme{default}

\usepackage{amsmath}
\usepackage{amssymb}
\usepackage{booktabs}

\setbeamersize{text margin left=0.55cm, text margin right=0.55cm}
\setlength{\parskip}{0.03cm}
\setbeamerfont{frametitle}{size=\large}
\setbeamertemplate{navigation symbols}{}

\title{Organic Chemistry Quiz Preparation}
\subtitle{Lesson $\rightarrow$ Example $\rightarrow$ Actual Quiz Question}
\author{}
\date{}

\begin{document}

%================================================
\begin{frame}
    \titlepage
\end{frame}

%================================================
\begin{frame}[t]{Topic 1: Oxidation State}
\small

\textbf{Lesson}

Oxidation state is found by assigning bonding electrons to the more electronegative atom.

\begin{itemize}
    \setlength{\itemsep}{0.06cm}
    \item Oxygen is usually $-2$.
    \item A neutral molecule has a total oxidation state of $0$.
    \item The more electronegative atom is treated as if it owns the bonding electrons.
\end{itemize}

\vspace{0.1cm}

\textbf{Example}

For carbon monoxide, CO:

\[
x+(-2)=0
\]

\[
x=+2
\]

So carbon has oxidation state:

\[
\boxed{+2}
\]

\end{frame}

%================================================
\begin{frame}[t]{Topic 1: Actual Quiz Question}
\scriptsize

\textbf{Actual Question}

What is the oxidation state and the formal charge on the carbon atom in carbon monoxide, respectively?

\vspace{0.08cm}

\begin{enumerate}
    \setlength{\itemsep}{0.04cm}
    \item[A.] $0;\ +1$
    \item[B.] $+1;\ 0$
    \item[C.] $+2;\ -1$
    \item[D.] $+4;\ 0$
\end{enumerate}

\vspace{0.08cm}

\textbf{Answer}

\[
\boxed{\text{C. } +2;\ -1}
\]

\vspace{-0.12cm}

\textbf{Why?}

Carbon has oxidation state $+2$ in CO.

Its formal charge in $:C\equiv O:$ is $-1$.

\end{frame}

%================================================
\begin{frame}[t]{Topic 2: Formal Charge}
\small

\textbf{Lesson}

Formal charge compares the electrons assigned to an atom with the number of valence electrons in the neutral atom.

\[
\text{Formal Charge}
=
\text{Valence}
-
\text{Nonbonding}
-
\frac{\text{Bonding}}{2}
\]

For carbon:

\[
\text{FC}
=
4-\text{Nonbonding}-\frac{\text{Bonding}}{2}
\]

\textbf{Important}

Formal charge and oxidation state are not the same thing.

\end{frame}

%================================================
\begin{frame}[t]{Topic 2: Example}
\small

\textbf{Example: Carbon in CO}

The Lewis structure of CO contains a triple bond:

\[
:C\equiv O:
\]

For carbon:

\begin{itemize}
    \setlength{\itemsep}{0.06cm}
    \item Valence electrons = $4$
    \item Nonbonding electrons = $2$
    \item Bonding electrons = $6$
\end{itemize}

\[
\text{FC}=4-2-\frac{6}{2}=4-2-3=-1
\]

\[
\boxed{\text{Formal Charge}=-1}
\]

\end{frame}

%================================================
\begin{frame}[t]{Topic 3: Oxygen in Carbon Dioxide}
\small

\textbf{Lesson}

The Lewis structure of carbon dioxide is:

\[
O=C=O
\]

Oxygen normally has oxidation state:

\[
\boxed{-2}
\]

Each oxygen has:

\begin{itemize}
    \setlength{\itemsep}{0.06cm}
    \item 6 valence electrons
    \item 4 nonbonding electrons
    \item 4 bonding electrons
\end{itemize}

\end{frame}

%================================================
\begin{frame}[t]{Topic 3: Example}
\small

\textbf{Example: Formal Charge of Oxygen}

For oxygen in CO$_2$:

\[
\text{FC}
=
6-4-\frac{4}{2}
\]

\[
\text{FC}=6-4-2=0
\]

Therefore:

\[
\boxed{\text{Oxidation State}=-2}
\qquad
\boxed{\text{Formal Charge}=0}
\]

\end{frame}

%================================================
\begin{frame}[t]{Topic 3: Actual Quiz Question}
\scriptsize

\textbf{Actual Question}

What is the oxidation state and the formal charge on the oxygen atom in carbon dioxide, respectively?

\vspace{0.08cm}

\begin{enumerate}
    \setlength{\itemsep}{0.04cm}
    \item[A.] $-2;\ 0$
    \item[B.] $-4;\ 0$
    \item[C.] $-1;\ +1$
    \item[D.] $+1;\ -2$
\end{enumerate}

\vspace{0.08cm}

\textbf{Answer}

\[
\boxed{\text{A. } -2;\ 0}
\]

\vspace{-0.12cm}

\textbf{Why?}

Oxygen normally has oxidation state $-2$.

Each oxygen in $O=C=O$ has formal charge $0$.

\end{frame}

%================================================
\begin{frame}[t]{Topic 4: Electronegativity and Bond Polarity}
\small

\textbf{Lesson}

Electronegativity describes how strongly an atom attracts bonding electrons.

\[
F > Cl > C > H > Li
\]

If carbon is bonded to an atom less electronegative than carbon, carbon pulls electrons toward itself.

\[
C^{\delta-}
\]

If carbon is bonded to an atom more electronegative than carbon:

\[
C^{\delta+}
\]

\end{frame}

%================================================
\begin{frame}[t]{Topic 4: Example}
\small

\textbf{Example: C--Li}

Carbon is more electronegative than lithium.

\[
C^{\delta-}-Li^{\delta+}
\]

So the carbon atom has a:

\[
\boxed{\text{partial negative charge}}
\]

Compare with C--F:

\[
C^{\delta+}-F^{\delta-}
\]

because fluorine is more electronegative than carbon.

\end{frame}

%================================================
\begin{frame}[t]{Topic 4: Actual Quiz Question}
\scriptsize

\textbf{Actual Question}

For each type of bond below containing a carbon, which carbon bears a partially negative charge?

\vspace{0.08cm}

\begin{enumerate}
    \setlength{\itemsep}{0.04cm}
    \item[A.] C--Li
    \item[B.] C--Cl
    \item[C.] C--C
    \item[D.] C--F
\end{enumerate}

\vspace{0.08cm}

\textbf{Answer}

\[
\boxed{\text{A. C--Li}}
\]

\vspace{-0.12cm}

\textbf{Why?}

Carbon is more electronegative than lithium.

Therefore carbon carries the partial negative charge.

\end{frame}

%================================================
\begin{frame}[t]{Topic 5: Skeletal Structures}
\small

\textbf{Lesson}

In skeletal or line-angle structures:

\begin{itemize}
    \setlength{\itemsep}{0.06cm}
    \item Every corner usually represents a carbon atom.
    \item Every unlabeled line ending usually represents a carbon.
    \item Hydrogens attached to carbon are normally omitted.
    \item Carbon normally forms a total of four bonds.
\end{itemize}

\[
\text{Hydrogens}
=
4-\text{bond order already used}
\]

\end{frame}

%================================================
\begin{frame}[t]{Topic 5: Examples}
\small

\textbf{Example 1: One Single Bond}

\[
4-1=3
\]

\[
\boxed{CH_3}
\]

\textbf{Example 2: Two Single Bonds}

\[
4-2=2
\]

\[
\boxed{CH_2}
\]

\textbf{Example 3: Double Bond + Single Bond}

\[
4-(2+1)=1
\]

\[
\boxed{CH}
\]

\end{frame}

%================================================
\begin{frame}[t]{Topic 5: Useful Groups}
\small

\begin{center}
\renewcommand{\arraystretch}{1.25}
\begin{tabular}{lc}
\toprule
\textbf{Group} & \textbf{Hydrogens} \\
\midrule
$-CH_3$ & 3 \\
$-CH_2-$ & 2 \\
$-CH=$ & 1 \\
$=CH_2$ & 2 \\
Internal alkyne carbon & 0 \\
Substituted aromatic carbon & 0 \\
Unsubstituted aromatic carbon & 1 \\
\bottomrule
\end{tabular}
\end{center}

\end{frame}

%================================================
\begin{frame}[t]{Topic 5: Actual Quiz Question}
\scriptsize

\textbf{Actual Question}

In molecular formulae, hydrogen atoms are often omitted for clarity.

How many hydrogens are present in the molecule shown in the quiz?

\vspace{0.08cm}

\begin{enumerate}
    \setlength{\itemsep}{0.04cm}
    \item[A.] 15
    \item[B.] 16
    \item[C.] 17
    \item[D.] 18
\end{enumerate}

\vspace{0.08cm}

\textbf{Answer}

\[
\boxed{\text{D. }18}
\]

\vspace{-0.12cm}

\textbf{Why?}

Count the hidden hydrogens on each carbon.

Each carbon must reach a total bond order of four.

\end{frame}

%================================================
\begin{frame}[t]{Topic 6: Hybridization}
\small

\textbf{Lesson}

Carbon uses different hybrid orbitals depending on its bonding environment.

\begin{center}
\renewcommand{\arraystretch}{1.25}
\begin{tabular}{lll}
\toprule
\textbf{Environment} & \textbf{Hybridization} & \textbf{Geometry} \\
\midrule
Single bonds & $sp^3$ & Tetrahedral \\
Double bond & $sp^2$ & Trigonal planar \\
Triple bond & $sp$ & Linear \\
\bottomrule
\end{tabular}
\end{center}

\[
\boxed{
\text{single} \rightarrow sp^3,\quad
\text{double} \rightarrow sp^2,\quad
\text{triple} \rightarrow sp
}
\]

\end{frame}

%================================================
\begin{frame}[t]{Topic 6: Examples}
\small

\textbf{Single Bonds}

\[
CH_4 \rightarrow \boxed{sp^3}
\]

\textbf{Double Bond}

\[
CH_2=CH_2 \rightarrow \boxed{sp^2}
\]

\textbf{Triple Bond}

\[
HC\equiv CH \rightarrow \boxed{sp}
\]

\textbf{Aromatic Ring}

\[
\text{benzene carbon} \rightarrow \boxed{sp^2}
\]

\end{frame}

%================================================
\begin{frame}[t]{Topic 7: Sigma Bonds}
\small

\textbf{Lesson}

A sigma bond, $\sigma$, forms through direct overlap between orbitals.

\[
\begin{array}{cc}
sp^3-sp^3 & sp^2-sp^2 \\[0.2cm]
sp-sp     & sp^2-sp
\end{array}
\]

\textbf{To determine orbital overlap:}

\begin{enumerate}
    \setlength{\itemsep}{0.06cm}
    \item Identify the two atoms forming the bond.
    \item Determine the hybridization of the first atom.
    \item Determine the hybridization of the second atom.
    \item Combine the two hybridizations.
\end{enumerate}

\end{frame}

%================================================
\begin{frame}[t]{Topic 7: Example}
\small

\textbf{Example}

Suppose a sigma bond connects:

\begin{itemize}
    \setlength{\itemsep}{0.06cm}
    \item a carbon in a benzene ring
    \item a carbon in a triple bond
\end{itemize}

The benzene carbon is:

\[
\boxed{sp^2}
\]

The alkyne carbon is:

\[
\boxed{sp}
\]

Therefore the sigma bond is:

\[
\boxed{sp^2-sp}
\]

\end{frame}

%================================================
\begin{frame}[t]{Topic 7: Actual Quiz Question}
\scriptsize

\textbf{Actual Question}

The indicated $\sigma$ bond results from the overlap of which orbitals?

\vspace{0.08cm}

\begin{enumerate}
    \setlength{\itemsep}{0.04cm}
    \item[A.] $sp^2-sp^2$
    \item[B.] $sp-sp^3$
    \item[C.] $sp^2-sp$
    \item[D.] $sp^2-sp^3$
\end{enumerate}

\vspace{0.08cm}

\textbf{Answer}

\[
\boxed{\text{C. }sp^2-sp}
\]

\vspace{-0.12cm}

\textbf{Why?}

The benzene-ring carbon is $sp^2$.

The adjacent triple-bond carbon is $sp$.

\end{frame}

%================================================
\begin{frame}[t]{Quiz Answer Key}
\small

\begin{center}
\renewcommand{\arraystretch}{1.35}
\begin{tabular}{ccc}
\toprule
\textbf{Question} & \textbf{Answer} & \textbf{Key Idea} \\
\midrule
1 & \textbf{C} & $+2;\,-1$ \\
2 & \textbf{A} & $-2;\,0$ \\
3 & \textbf{A} & C--Li \\
4 & \textbf{D} & 18 hydrogens \\
5 & \textbf{C} & $sp^2-sp$ \\
\bottomrule
\end{tabular}
\end{center}

\vspace{0.2cm}

\begin{center}
\textbf{Understand the reason, not only the letter.}
\end{center}

\end{frame}

%================================================
\begin{frame}[t]{Quick Review}
\small

For this quiz, remember:

\begin{enumerate}
    \setlength{\itemsep}{0.08cm}
    \item Oxidation state
    \item Formal charge
    \item Electronegativity
    \item Bond polarity
    \item Skeletal structures
    \item Counting hidden hydrogens
    \item Carbon hybridization
    \item Sigma-bond orbital overlap
\end{enumerate}

\end{frame}

%================================================
\begin{frame}[t]{Problem-Solving Strategy}
\small

When answering the quiz:

\begin{enumerate}
    \setlength{\itemsep}{0.08cm}
    \item Draw or inspect the Lewis structure.
    \item Count bonds and lone pairs.
    \item Calculate formal charge.
    \item Use electronegativity for oxidation state and polarity.
    \item Complete carbon's four bonds when counting hydrogens.
    \item Look for single, double, or triple bonds.
    \item Determine the hybridization of each relevant atom.
    \item Determine the orbital overlap.
\end{enumerate}

\end{frame}

\end{document}
